\subsection{The Most Repeatable Form Was Diffuse Tangential Activity}

\begin{figure*}[t]
\centering
\includegraphics[width=\textwidth]{figures/Fig3_final.pdf}
\caption{Spatial and timing structure of the T1 residual. Support fractions
and aggregate event-aligned profiles show that diffuse all-tangential activity
was the most stable repeated form. Near-pre timing, edge/core contrast, and
signed structure were retained as bounded or secondary patterns.}
\label{fig:spatial_timing}
\end{figure*}

After establishing survival and local robustness, we asked what form the
surviving residual took across space and event time. The most stable repeated
form was diffuse tangential activity: tangential non-affine activity
distributed across sampled focal neighborhoods rather than concentrated in a
stable edge/core contrast, a fixed near-pre event window, or a universal
signed transition direction.

Diffuse all-tangential activity was supported in 13 of the 14 tested
observations, whereas near-pre timing was supported in 8 of the 14 tested
observations (Fig.~\ref{fig:spatial_timing}). Edge/core contrast appeared in a
subset of observations, but it was weaker as a phase-specific interpretation.
Signed event-type structure also remained heterogeneous and did not define a
stable low-to-high, high-to-low, or mirror-symmetric rule. The supported
phenotype is therefore distributed local tangential non-affine activity, not a
sharp event precursor, stable edge trigger, or global signed law.

The near-pre enrichment should be read with this boundary in mind. It was a
descriptive pattern under the original event-control comparison and did not
survive the later state-matched event-locality challenge described below. The
detrending-specific audit of the same near-pre window retained majority-level
support under centered, past-only, and no-rolling variants, but that audit used
a looser phase-tail gate and a different pseudo-event sampler. We therefore
treat those counts as sensitivity evidence rather than as replacements for the
original $8/14$ phase-localization count. Thus, the near-pre profile remained a
moderate timing description, not evidence for a causal online precursor.
